\documentclass[11pt,a4paper]{article}
\usepackage[margin=25mm]{geometry}
\usepackage{amsmath,amssymb,amsthm,booktabs,hyperref}
\setlength{\emergencystretch}{2em}
\newtheorem{theorem}{Theorem}
\newtheorem{lemma}[theorem]{Lemma}
\newcommand{\PSL}{\operatorname{PSL}}
\newcommand{\SL}{\operatorname{SL}}
\newcommand{\Aut}{\operatorname{Aut}}
\title{Two exact columns in the O'Nan map enumeration}
\author{Human Compression project}
\date{}
\begin{document}\maketitle
\begin{abstract}
The maximal-subgroup bound becomes an equality whenever each relevant
nongenerating pair has a unique maximal overgroup. We prove this for
all rotation types involving 19 or 31 in the O'Nan group. For order 19,
the two classes of $L_3(7):2$ are essential, as well as $J_1$.
The argument uses the established maximal-subgroup classifications and
the previously acquired character vectors. It determines 31 distinct
table entries, of which 21 were not determined by the preceding bounds.
The other entries and the chirality spectrum remain separate questions.
\end{abstract}

\section{An exact subtraction principle}
For a finite group $H$ put
\[
 X_H(m,p)=\{(r,s)\in H^2: |r|=m,\ |s|=p,\ |rs|=2\},
 \qquad x_H(m,p)=|X_H(m,p)|.
\]
Let $E_H(m,p)$ be the subset generating $H$, and write $e_H=|E_H|$.
These are ordered pairs. Automorphisms act freely on generating pairs,
so the number of generating-pair orbits is $e_H/|\Aut(H)|$.
This is the convention of formula (2.4) and Table 1 in
Connor--Leemans~\cite{CL}; there is no further division by duality.

\begin{lemma}\label{unique}
Let $G$ be a nonabelian finite simple group. Suppose every pair in
$X_G(m,p)\setminus E_G(m,p)$ lies in exactly one maximal proper subgroup.
If $M_1,\ldots,M_t$ represent the conjugacy classes of maximal subgroups
whose orders are divisible by $p$, then
\[
 e_G(m,p)=x_G(m,p)-\sum_{j=1}^t[G:M_j]x_{M_j}(m,p). \tag{1}
\]
\end{lemma}
\begin{proof}
Every nongenerating pair lies in a maximal proper subgroup, and the
order of that subgroup is divisible by $p$. Conversely, every pair in
any such maximal subgroup fails to generate $G$. The uniqueness
hypothesis makes these sets of pairs disjoint. Every maximal proper
subgroup $M$ is self-normalizing: maximality gives $N_G(M)=M$ or $G$,
and the second case would make the nontrivial proper subgroup $M$
normal in $G$. Thus its conjugacy class has size $[G:M]$. Summing the
sizes of the disjoint sets proves (1). Distinct conjugacy classes of
isomorphic maximal subgroups are counted separately.
\end{proof}

\begin{lemma}\label{normalizer}
Let $G$ be simple, let $M$ be maximal, and let $1<H\lhd M$ be proper
in $G$. Suppose every maximal subgroup of $G$ containing $H$ has
$H$ as a subgroup of index two. Then $M$ is the unique maximal
subgroup containing $H$.
\end{lemma}
\begin{proof}
The normalizer of $H$ contains $M$. It cannot be $G$, since $H$ is
nontrivial and proper. Hence $N_G(H)=M$. Any maximal overgroup $K$ of
$H$ normalizes $H$ because $[K:H]=2$, and therefore $K\leq M$.
Maximality implies $K=M$.
\end{proof}

\section{The column involving 31}
We first prove the required fact about $\PSL_2(p)$ directly.

\begin{lemma}\label{prime}
Let $p$ be an odd prime and $L=\PSL_2(\mathbb F_p)$ in its natural
action on the projective line. Every proper subgroup of $L$ containing
an element of order $p$ is contained in a point stabilizer.
\end{lemma}
\begin{proof}
An element of order $p$ lifts, after changing the sign if necessary,
to a nonidentity unipotent matrix in $\SL_2(\mathbb F_p)$.
Indeed its eigenvalue ratio has $p$th power one in characteristic $p$,
so both eigenvalues coincide, and determinant one makes them $1$ or
$-1$. A nonidentity unipotent matrix has exactly one projective fixed
point. The point stabilizer at $\infty$ consists of upper triangular
matrices modulo the scalar centre. Its unique Sylow $p$ subgroup is
\[
 U_\infty=\left\{\begin{bmatrix}1&a\\0&1\end{bmatrix}:a\in\mathbb F_p\right\}
 \quad\text{modulo the centre}.
\]
An order-$p$ subgroup fixing $\infty$ is therefore all of $U_\infty$.
Here the assumption that the field has prime order is essential.

Let $K$ contain $U_\infty$. If $K$ does not stabilize $\infty$, it
contains a conjugate $U_y$ with $y\ne\infty$. Translation by an element
of $U_\infty$ takes $y$ to $0$, so $K$ contains $U_0$ as well.
The corresponding upper and lower elementary matrices generate
$\SL_2(\mathbb F_p)$. To see this explicitly, elementary row operations
reduce any determinant-one matrix to a diagonal matrix. The remaining
diagonal matrices also lie in the generated group: if
$u(a)=\left[\begin{smallmatrix}1&a\\0&1\end{smallmatrix}\right]$
and $v(b)=\left[\begin{smallmatrix}1&0\\b&1\end{smallmatrix}\right]$,
then
\[
 w(a)=u(a)v(-a^{-1})u(a)=\begin{bmatrix}0&a\\-a^{-1}&0\end{bmatrix},
 \qquad w(a)w(-1)=\begin{bmatrix}a&0\\0&a^{-1}\end{bmatrix}.
\]
Thus $K=L$, a contradiction. This proves the lemma.
\end{proof}

In $\PSL_2(31)$ the point stabilizer has order
$31(31-1)/2=465$, which is odd. Consequently a proper subgroup
containing an element of order 31 cannot contain an involution.
Every pair in $X_{\PSL_2(31)}(m,31)$ therefore generates
$\PSL_2(31)$, irrespective of $m$.

The complete maximal-subgroup list of the O'Nan group $G$, acquired in
the preceding bounds component from the Atlas~\cite{AtlasON}, contains
exactly two classes with order divisible by 31. Both are isomorphic
to $L_2(31)$ and have order 14880. Any nongenerating pair of the
prescribed type generates its containing maximal subgroup. That
subgroup is consequently unique. Formula (1) gives
\[
 e_G(m,31)=x_G(m,31)-2[G:L_2(31)]x_{L_2(31)}(m,31). \tag{2}
\]
The actual certificate retains both table identifiers, \texttt{L2(31)}
and \texttt{ONM8}, and adds their separately evaluated contributions.
It checks their equality without discarding either conjugacy class.

\section{The column involving 19}
The order-19 overgroups require a different argument. In particular,
$|L_3(7):2|=3753792$ is divisible by 19. Omitting those maximal
subgroups would give an incorrect subtraction.

\begin{lemma}\label{cyclic}
Suppose $K$ has a normal subgroup $P$ of odd prime order $p$, and
$K$ contains no element of order $2p$. If $r\in K$, $s\in P$ has order $p$
and $rs$ is an involution, then $r$ is an involution.
\end{lemma}
\begin{proof}
Put $b=rs$. Conjugation by $b$ induces an automorphism of $P$ whose
order divides two. It cannot be trivial: otherwise $b$ and $s$ commute,
and their product has order $2p$. The unique automorphism of order two
of the cyclic group $P$ is inversion. Hence
$(bs^{-1})^2=1$. Since $b\notin P$, the element $r=bs^{-1}$ is not
the identity and has order two.
\end{proof}

We use the complete maximal-subgroup lists in the primary Atlas pages
for $J_1$, $L_3(7)$ and $L_3(7):2$~\cite{AtlasJ1,AtlasL37}.
Their relevant consequences are as follows.
\begin{itemize}
\item In $J_1$, the only maximal subgroups with order divisible by 19
are of type $19:6$.
\item In $L=L_3(7)$, the only such maximal subgroups have type $19:3$.
\item In $M=L_3(7):2$, they are the normal subgroup $L$ and subgroups
of type $19:6$.
\end{itemize}
The groups of type $19:6$ have a normal subgroup of order 19, containing
every element of that order. The acquired character tables show that
neither $J_1$ nor $M$ has an element of order 38. Lemma~\ref{cyclic}
therefore excludes all pairs of our type with $m>2$ from these $19:6$
subgroups. The groups $19:3$ have odd order, so contain no involution.

It follows that every pair in $X_{J_1}(m,19)$ with $m>2$ generates
$J_1$. A pair in $X_M(m,19)$ either generates $M$ or is contained in
$L$; in the latter case it generates $L$, by the list of maximal
subgroups of $L$ and the odd-order observation.

Among the thirteen maximal classes of $G$, exactly three have order
divisible by 19: one class of $J_1$ and two classes of $M$.
A pair generating one of these maximal subgroups has a unique maximal
overgroup. It remains to consider a pair generating $L$.
The complete list and subgroup orders show that every maximal overgroup
of $L$ in $G$ is isomorphic to $M$, with $L$ of index two.
Lemma~\ref{normalizer} gives uniqueness in this case also. Therefore
\[
 e_G(m,19)=x_G(m,19)-[G:J_1]x_{J_1}(m,19)
                    -2[G:M]x_M(m,19),\qquad m>2. \tag{3}
\]
Again the two $M$ contributions are retained separately under the names
\texttt{L3(7).2} and \texttt{ONM2}. No assertion about pairs with $m=2$
is included in (3).

\section{Exact evaluation and source comparison}
The acquired common counting theorem gives, for any retained table $H$,
\[
 x_H(m,p)=\frac{1}{|H|}\sum_{\chi\in\operatorname{Irr}(H)}
 \frac{v_{\chi,2}v_{\chi,m}v_{\chi,p}}{\chi(1)},\qquad
 v_{\chi,t}=\sum_{|c_i|=t}|C_i|\chi(c_i). \tag{4}
\]
The saved vectors are integral. Our reader clears the character-degree
denominators by their least common multiple and performs only exact
integer operations. The maximal terms in (2) and (3) count raw ordered
pairs; division by $|\Aut(G)|=921631011840$ is performed afterward.
No normalized count for a smaller group is substituted for a raw count.
All displayed results have zero remainder under this division.

\begin{center}\begin{tabular}{r|rr}\toprule
$m$&$n_G(m,19)$&$n_G(m,31)$\\\midrule
3&57&25\\
4&846&359\\
5&1242&502\\
6&2943&1268\\
7&4725&1995\\
8&13926&5752\\
10&11262&4632\\
11&20493&8325\\
12&6126&2577\\
14&7899&3266\\
15&10020&4062\\
16&56052&23048\\
19&35442&14553\\
20&22572&9118\\
28&15978&6658\\
31&14553&5999\\
\bottomrule\end{tabular}\end{center}


Swapping the two elements is a bijection of the corresponding type
sets, since $sr$ is conjugate to $rs$. Thus the two displayed columns
treat 31 distinct source cells: the cell $(19,31)$ appears in both and
is counted once. Ten of these cells were already exact in the preceding
bounds. The new component determines 21 more, for a combined total of
41 exact cells among 136. It preserves the correction $9118$, rather
than $7246$, for $(20,31)$; the other 30 displayed cells agree with the
fixed source table. These comparisons are made only after evaluation.

The unique-overgroup argument and its explicit application are the
contribution of this component. The O'Nan maximal classification,
its automorphism order, the smaller maximal-subgroup lists and all
ordinary character tables are borrowed and attributed. We do not
claim a new classification, a conjecture settlement, a full source
recovery, Lean formalization or a measured compression factor.
The 95 remaining exact table entries and the source's inversion and
chirality assertions remain unpaid by these components.

\begin{thebibliography}{9}
\bibitem{CL} T. Connor and D. Leemans, \emph{Algorithmic enumeration of
regular maps}, Ars Mathematica Contemporanea 10 (2016), 211--222.
\bibitem{Wilson} R. A. Wilson, \emph{The maximal subgroups of the O'Nan
group}, Journal of Algebra 97 (1985), 467--473.
\bibitem{AtlasON} R. A. Wilson, S. J. Nickerson, J. N. Bray et al.,
\emph{Atlas of Group Representations}, O'Nan group page,
\url{https://brauer.maths.qmul.ac.uk/Atlas/v3/spor/ON/}.
\bibitem{AtlasJ1} The same authors, Janko group $J_1$ page,
\url{https://brauer.maths.qmul.ac.uk/Atlas/v3/spor/J1/},
maximal-subgroup section, accessed 16 October 2026.
\bibitem{AtlasL37} The same authors, $L_3(7)$ page,
\url{https://brauer.maths.qmul.ac.uk/Atlas/v3/lin/L37/},
maximal-subgroup sections for $L_3(7)$ and $L_3(7):2$,
accessed 16 October 2026.
\end{thebibliography}
\end{document}
